% ============================================================
%  Análisis de Puntos de Función — Sistema de Gestión de Notas
%  Universidad de las Fuerzas Armadas "ESPE"
%  Para compilar en Overleaf: pdflatex
% ============================================================

\documentclass[12pt, a4paper]{article}

% ── Paquetes ─────────────────────────────────────────────────
\usepackage[utf8]{inputenc}
\usepackage[T1]{fontenc}
\usepackage[spanish]{babel}
\usepackage{geometry}
\usepackage{graphicx}
\usepackage{booktabs}
\usepackage{tabularx}
\usepackage{longtable}
\usepackage{array}
\usepackage{xcolor}
\usepackage{colortbl}
\usepackage{hyperref}
\usepackage{setspace}
\usepackage{parskip}
\usepackage{titlesec}
\usepackage{fancyhdr}
\usepackage{lmodern}

% ── Márgenes ─────────────────────────────────────────────────
\geometry{
  top    = 2.5cm,
  bottom = 2.5cm,
  left   = 3cm,
  right  = 2.5cm
}

% ── Colores institucionales ───────────────────────────────────
\definecolor{espeazul}{RGB}{0, 51, 102}
\definecolor{grisclaro}{RGB}{245, 245, 245}
\definecolor{grisoscuro}{RGB}{80, 80, 80}
\definecolor{verdeok}{RGB}{0, 128, 0}

% ── Estilo de secciones ───────────────────────────────────────
\titleformat{\section}
  {\large\bfseries\color{espeazul}}
  {\thesection.}{0.5em}{}[\titlerule]

\titleformat{\subsection}
  {\normalsize\bfseries\color{grisoscuro}}
  {\thesubsection.}{0.5em}{}

% ── Encabezado y pie de página ────────────────────────────────
\pagestyle{fancy}
\fancyhf{}
\fancyhead[L]{\small\color{grisoscuro} Análisis de Puntos de Función — Gestión de Notas}
\fancyhead[R]{\small\color{grisoscuro} ESPE · NRC 30744}
\fancyfoot[C]{\small\thepage}
\renewcommand{\headrulewidth}{0.4pt}

% ── Colores de filas de tabla ─────────────────────────────────
\rowcolors{2}{grisclaro}{white}

% =============================================================
\begin{document}

% ── CARÁTULA ─────────────────────────────────────────────────
\begin{titlepage}
  \centering
  \vspace*{1cm}

  \includegraphics[width=5cm]{img/logoESPE}

  \vspace{1.2cm}

  {\Large\bfseries\color{espeazul}
    UNIVERSIDAD DE LAS FUERZAS ARMADAS ``ESPE''
  }

  \vspace{0.4cm}
  {\large\bfseries\color{grisoscuro}
    DEPARTAMENTO DE CIENCIAS DE LA COMPUTACIÓN
  }

  \vspace{2cm}
  \rule{0.85\textwidth}{0.6pt}
  \vspace{0.5cm}

  {\LARGE\bfseries\color{espeazul}
    Análisis de Puntos de Función\\[0.4em]
    \large Sistema Web de Gestión de Notas
  }

  \vspace{0.5cm}
  \rule{0.85\textwidth}{0.6pt}

  \vspace{2cm}

  \begin{tabular}{rl}
    \textbf{Nombre:}  & Luis Cueva                              \\[6pt]
    \textbf{NRC:}     & 30744                                   \\[6pt]
    \textbf{Materia:} & Construcción y Evolución del Software   \\[6pt]
    \textbf{Fecha:}   & 4 de mayo de 2026                       \\
  \end{tabular}

  \vfill
  {\small\color{grisoscuro} Sangolquí, Ecuador}
\end{titlepage}

% ── CUERPO DEL DOCUMENTO ─────────────────────────────────────
\newpage
\setcounter{page}{1}
\onehalfspacing

% =============================================================
\section{Introducción}

El presente documento detalla la estimación del tamaño funcional del sistema web
\textbf{Gestión de Notas}, una plataforma genérica desarrollada para
\textbf{instituciones educativas de nivel básico y bachillerato},
mediante la técnica de \textbf{Puntos de Función (Function Points)} bajo el
estándar \textbf{IFPUG} (\textit{International Function Point Users Group}).

\vspace{0.3cm}
\begin{tabular}{ll}
  \textbf{Stack tecnológico:} & Next.js · NestJS · PostgreSQL · Prisma ORM \\
  \textbf{Método aplicado:}   & IFPUG — estimación por componente funcional \\
  \textbf{Roles del sistema:} & Administrador, Profesor, Estudiante          \\
\end{tabular}

% =============================================================
\section{Referencia de Pesos por Tipo y Complejidad (IFPUG)}

\begin{table}[h!]
  \centering
  \caption{Tabla de pesos estándar IFPUG por tipo de componente funcional}
  \vspace{0.3cm}
  \begin{tabular}{>{\bfseries}p{3.5cm} p{6cm} >{\centering\arraybackslash}p{1.5cm} >{\centering\arraybackslash}p{1.5cm} >{\centering\arraybackslash}p{1.5cm}}
    \toprule
    \rowcolor{espeazul}
    \textcolor{white}{Tipo} &
    \textcolor{white}{Descripción} &
    \textcolor{white}{Baja} &
    \textcolor{white}{Media} &
    \textcolor{white}{Alta} \\
    \midrule
    EI — Entradas Externas    & Datos que ingresan al sistema (formularios, uploads) & 3 & 4 & 6  \\
    EO — Salidas Externas     & Datos derivados que salen del sistema (reportes, PDFs, cálculos) & 4 & 5 & 7  \\
    EQ — Consultas Externas   & Consultas sin transformación (listados, búsquedas)   & 3 & 4 & 6  \\
    ILF — Arch. Lógicos Int.  & Entidades persistidas en la base de datos            & 7 & 10 & 15 \\
    \bottomrule
  \end{tabular}
\end{table}

% =============================================================
\section{Tabla de Puntos de Función por Historia de Usuario}

\begin{center}
\scriptsize
\begin{longtable}{>{\bfseries}p{1cm} p{5.8cm} >{\centering\arraybackslash}p{0.6cm} >{\centering\arraybackslash}p{0.6cm} >{\centering\arraybackslash}p{0.6cm} >{\centering\arraybackslash}p{0.6cm} >{\centering\arraybackslash}p{1.4cm} >{\centering\arraybackslash}p{1.2cm}}

  \toprule
  \rowcolor{espeazul}
  \textcolor{white}{\textbf{ID}} &
  \textcolor{white}{\textbf{Historia de Usuario}} &
  \textcolor{white}{\textbf{EI}} &
  \textcolor{white}{\textbf{EO}} &
  \textcolor{white}{\textbf{EQ}} &
  \textcolor{white}{\textbf{ILF}} &
  \textcolor{white}{\textbf{Complejidad}} &
  \textcolor{white}{\textbf{PF Est.}} \\
  \midrule
  \endfirsthead

  \toprule
  \rowcolor{espeazul}
  \textcolor{white}{\textbf{ID}} &
  \textcolor{white}{\textbf{Historia de Usuario}} &
  \textcolor{white}{\textbf{EI}} &
  \textcolor{white}{\textbf{EO}} &
  \textcolor{white}{\textbf{EQ}} &
  \textcolor{white}{\textbf{ILF}} &
  \textcolor{white}{\textbf{Complejidad}} &
  \textcolor{white}{\textbf{PF Est.}} \\
  \midrule
  \endhead

  \bottomrule
  \endfoot

  HU-01 & Como usuario quiero iniciar sesión con email y contraseña para acceder al sistema con mi rol correspondiente.
         & 1 & 1 & 1 & 1 & Baja & \textbf{14} \\
  HU-02 & Como usuario quiero cerrar sesión para invalidar mi token y proteger el acceso en equipos compartidos.
         & 1 & 0 & 1 & 0 & Baja & \textbf{6} \\
  HU-03 & Como administrador quiero crear, editar y desactivar cuentas de usuario con rol asignado (Admin, Profesor, Estudiante).
         & 3 & 1 & 1 & 1 & Media & \textbf{24} \\
  HU-04 & Como administrador quiero crear y gestionar años lectivos, marcando uno como activo para el período operativo actual.
         & 2 & 1 & 1 & 1 & Baja & \textbf{20} \\
  HU-05 & Como administrador quiero crear cursos, crear materias y asignar materias a cursos dentro de un año lectivo.
         & 3 & 1 & 2 & 2 & Media & \textbf{36} \\
  HU-06 & Como administrador quiero matricular estudiantes en cursos y trasladarlos entre secciones conservando el historial.
         & 2 & 1 & 2 & 2 & Media & \textbf{31} \\
  HU-07 & Como profesor quiero crear los 3 parciales de un curso-materia para estructurar el período de evaluación.
         & 1 & 1 & 1 & 1 & Baja & \textbf{17} \\
  HU-08 & Como profesor quiero crear actividades (TAREA, PRUEBA, PROYECTO, EXAMEN) con fechas y validaciones de límite por tipo por parcial.
         & 3 & 1 & 2 & 2 & Alta & \textbf{51} \\
  HU-09 & Como profesor quiero revisar el PDF de evidencia de cada estudiante y asignar una nota de 0 a 10 por actividad.
         & 2 & 2 & 2 & 2 & Alta & \textbf{48} \\
  HU-10 & Como estudiante quiero ver las actividades de cada parcial con indicadores visuales de estado de entrega (Entregada / Pendiente / No entregada).
         & 1 & 2 & 2 & 1 & Media & \textbf{27} \\
  HU-11 & Como estudiante quiero subir un archivo PDF como evidencia de una actividad (máximo 10 MB, dentro del plazo).
         & 2 & 1 & 1 & 2 & Media & \textbf{33} \\
  HU-12 & Como estudiante quiero reemplazar mi evidencia entregada mientras el plazo de la actividad esté vigente.
         & 2 & 1 & 1 & 1 & Baja & \textbf{20} \\
  HU-13 & Como sistema quiero calcular automáticamente el promedio ponderado de cada parcial (TAREA 20\%, PRUEBA 20\%, PROYECTO 25\%, EXAMEN 35\%) y el promedio final de la materia.
         & 2 & 3 & 2 & 3 & Alta & \textbf{66} \\
  HU-14 & Como sistema quiero aplicar la regla de supletorio cuando el promedio final sea menor a 7.00, actualizando el estado del estudiante automáticamente.
         & 2 & 2 & 2 & 2 & Alta & \textbf{48} \\
  HU-15 & Como sistema quiero enviar notificaciones automáticas al estudiante cuando se registre una calificación o cambie su estado de aprobación.
         & 1 & 2 & 1 & 2 & Media & \textbf{33} \\
  HU-16 & Como sistema quiero ejecutar un Cron Job periódico que elimine archivos BYTEA de evidencias antiguas para liberar espacio en base de datos.
         & 1 & 1 & 1 & 2 & Media & \textbf{27} \\

\end{longtable}
\end{center}

% =============================================================
\section{Totales por Módulo}

\begin{table}[h!]
  \centering
  \caption{Resumen de Puntos de Función por módulo del sistema}
  \vspace{0.3cm}
  \begin{tabular}{p{7.5cm} >{\centering\arraybackslash}p{3cm} >{\centering\arraybackslash}p{2cm}}
    \toprule
    \rowcolor{espeazul}
    \textcolor{white}{\textbf{Módulo}} &
    \textcolor{white}{\textbf{Historias}} &
    \textcolor{white}{\textbf{PF Total}} \\
    \midrule
    Autenticación y Control de Roles              & HU-01, HU-02                       & 20  \\
    CRUD Administrativo                           & HU-03, HU-04, HU-05, HU-06         & 111 \\
    Gestión Transaccional Docente                 & HU-07, HU-08, HU-09                & 116 \\
    Gestión Transaccional Estudiante              & HU-10, HU-11, HU-12                & 80  \\
    Motor de Automatización (Calificaciones)      & HU-13, HU-14, HU-15                & 147 \\
    Mantenimiento (Cron Job)                      & HU-16                              & 27  \\
    \midrule
    \rowcolor{espeazul}
    \textcolor{white}{\textbf{TOTAL DEL SISTEMA}} &
    \textcolor{white}{\textbf{16 historias}}       &
    \textcolor{white}{\textbf{501}}                \\
    \bottomrule
  \end{tabular}
\end{table}

% =============================================================
\section{Justificación Académica}

El presente análisis de Puntos de Función evidencia que el módulo de
\textbf{Calificaciones y Promedios} (HU-13 y HU-14) constituye el componente
de mayor complejidad funcional del sistema, acumulando un total de
\textbf{114 Puntos de Función} entre ambas historias. Esta complejidad se
origina en la densidad de sus reglas de negocio: el cálculo del promedio
parcial requiere una ponderación diferenciada por tipo de actividad
(TAREA, PRUEBA, PROYECTO y EXAMEN), mientras que el promedio final de materia
integra los resultados de tres parciales independientes y activa, de manera
condicional, la regla de supletorio cuando el resultado es inferior a 7.00.
La combinación de múltiples Archivos Lógicos Internos (ILF), Salidas
Externas (EO) de datos derivados y Consultas Externas (EQ) cruzadas entre los
módulos de actividades, calificaciones y promedios sitúa a este componente en
el nivel de \textbf{complejidad Alta} según el estándar IFPUG, justificando
el mayor esfuerzo de diseño, desarrollo y pruebas asignado durante el Sprint~5
del cronograma del proyecto.

\vfill
\begin{center}
  \small\color{grisoscuro}
  Sistema de Gestión de Notas --- Plataforma para Instituciones Educativas \quad|\quad
  Análisis IFPUG v1.0 \quad|\quad
  Universidad de las Fuerzas Armadas ``ESPE''
\end{center}

\end{document}
